\chapter{Estruturas de Decisão}
\label{cap:decisao}

Até aqui, todo código apresentado executa suas instruções sequencialmente, uma após a outra, do início ao fim. Estruturas de decisão permitem que o programa escolha \emph{qual} caminho seguir, dependendo de uma condição — é isso que dá a um programa a capacidade de reagir de forma diferente a diferentes situações.

\section{\code{if} / \code{else}}

\begin{codigoc}{Estrutura if / else if / else}
#include <stdio.h>

int main(void) {
    int temperatura = 32;

    if (temperatura > 30) {
        printf("Esta quente!\n");
    } else if (temperatura < 15) {
        printf("Esta frio!\n");
    } else {
        printf("Temperatura agradavel.\n");
    }

    return 0;
}
\end{codigoc}

Em C, \textbf{qualquer valor diferente de zero é considerado verdadeiro}; apenas o valor \code{0} é considerado falso. Isso significa que \code{if (10)} sempre executa seu bloco, e \code{if (0)} nunca executa — mesmo sem envolver nenhum operador relacional explícito.

\begin{dica}
Quando o bloco de um \code{if} (ou de qualquer estrutura de controle) contém uma única instrução, as chaves são opcionais. Ainda assim, é considerada boa prática sempre usá-las — isso evita um erro clássico onde uma segunda linha, adicionada depois por engano fora das chaves, passa a ser executada incondicionalmente.
\end{dica}

\section{Operador ternário}

Para condicionais simples que retornam um valor, C oferece uma forma compacta:

\begin{codigoc}{Operador ternário}
int a = 10, b = 20;
int maior = (a > b) ? a : b;   // "se a > b, entao a, senao b"

printf("O maior valor e %d\n", maior);
\end{codigoc}

\section{\code{switch}}

Quando uma mesma variável precisa ser comparada contra vários valores possíveis, o \code{switch} costuma ser mais legível do que uma longa cadeia de \code{else if}:

\begin{codigoc}{Estrutura switch}
#include <stdio.h>

int main(void) {
    int dia = 3;

    switch (dia) {
        case 1:
            printf("Segunda-feira\n");
            break;
        case 2:
            printf("Terca-feira\n");
            break;
        case 3:
            printf("Quarta-feira\n");
            break;
        default:
            printf("Dia invalido\n");
    }

    return 0;
}
\end{codigoc}

\begin{atencao}
Esquecer o \code{break} ao final de um \code{case} faz a execução ``cair'' (\emph{fall-through}) para o próximo \code{case}, executando-o também — mesmo que sua condição não tenha sido satisfeita. Às vezes esse comportamento é usado intencionalmente (para agrupar vários casos com a mesma ação), mas na grande maioria das vezes é um bug. O \code{default}, por convenção o último caso, captura qualquer valor não tratado pelos \code{case}s anteriores.
\end{atencao}

\begin{esp32box}
Estruturas \code{switch} são extremamente comuns em firmware para implementar \textbf{máquinas de estado} — por exemplo, controlar as diferentes fases de operação de um dispositivo (``aguardando'', ``conectando ao Wi-Fi'', ``lendo sensor'', ``enviando dados'', ``erro''), onde cada estado é representado por um valor de um \code{enum} (\cref{cap:structs}) e o \code{switch} decide o que fazer em cada um deles a cada iteração do laço principal do programa.
\end{esp32box}
